% Generated draft: requires mathematical and source review.
\documentclass[11pt]{article}
\usepackage[margin=25mm]{geometry}
\usepackage{fontspec}
\setmainfont{Latin Modern Roman}
\usepackage{amsmath,amssymb,mathtools,unicode-math}
\setmathfont{Latin Modern Math}
\usepackage{enumitem,needspace,xurl,hyperref}
\setlength{\parindent}{0pt}
\setlength{\parskip}{5pt}
\setlength{\emergencystretch}{3em}
\newcommand{\N}{\mathbb{N}}
\newcommand{\Z}{\mathbb{Z}}
\newcommand{\Q}{\mathbb{Q}}
\newcommand{\R}{\mathbb{R}}
\newcommand{\C}{\mathbb{C}}
\newcommand{\F}{\mathbb{F}}
\newcommand{\A}{\mathbb{A}}
\newcommand{\PP}{\mathbb{P}}
\newcommand{\E}{\mathbb{E}}
\newcommand{\eps}{\varepsilon}
\newcommand{\dd}{\,\mathrm{d}}
\begin{document}
\section*{UnsolvedMath ID: 701392}
\textbf{Source problem code:} catalogue entry
\subsection{Definitions and mathematical statement}
For a group \(G\), define its lower central series by
\[
\gamma_1(G)=G,\qquad \gamma_{n+1}(G)=[G,\gamma_n(G)]\quad(n\geq 1),
\]
where \([A,B]\) is the subgroup generated by the commutators \(a^{-1}b^{-1}ab\) with \(a\in A\) and \(b\in B\). The group \(G\) is residually nilpotent if
\[
\bigcap_{n\geq 1}\gamma_n(G)=\{1\}.
\]
A group \(G\) is called parafree if it is residually nilpotent and there exists a free group \(F\) such that
\[
G/\gamma_n(G)\cong F/\gamma_n(F)\qquad\text{for every integer }n\geq 1.
\]
Here the source's phrase ``the same lower central quotients'' is interpreted in the standard sense of these quotient groups, with one fixed free group \(F\).

\textbf{Question (Kourovka Notebook, Problem 17.90).} For every finitely generated parafree group \(G\), must
\[
H_2(G;\Z)=0
\]
hold? Here \(H_2(G;\Z)\) denotes second group homology with trivial integral coefficients. The source writes \(H_2(G)\) without specifying coefficients; the integral interpretation follows the usual convention and requires confirmation in review.
\subsection{Short English statement}
Must every finitely generated parafree group have trivial second integral homology? A parafree group is residually nilpotent and has the same lower central quotient groups as a free group. Integral coefficients are understood by convention because the source does not specify them.
\subsection{Sources}
\begin{itemize}
\item Local UnsolvedMath registry, supplied complete source record: permanent ID \(701392\), exact source code \texttt{catalogue entry}, title ``Kourovka Notebook Problem 17.90''; upstream revision \texttt{d371ee917ac9647fbb06cd6fc529deec6f3dda82}.
\item Supplied reference: \emph{The Kourovka Notebook (Unsolved Problems in Group Theory)}, 17th issue (2010), Problem 17.90; also identified in the record with the 21st issue edition (2026), arXiv:\texttt{1401.0300v48}. The collection description names E. I. Khukhro and V. D. Mazurov as editors. These bibliographic details are supplied attribution and have not been independently verified.
\item Attribution: the proposer fields and background name R. Mikhailov, with proposed year 2010; the statement itself begins ``(G. Baumslag).'' Both attributions are retained without resolving their respective roles.
\item The local record labels the problem \emph{open} and records an update on 2026-10-16. This is a supplied registry status, not an independently verified assessment of the literature.
\end{itemize}
\end{document}
